\documentclass[12pt,prd,preprintnumbers,amsmath,amssymb]{revtex4-2}
\usepackage{graphicx}
\usepackage{booktabs}
\usepackage{hyperref}
\usepackage{amsmath,amssymb,amsthm}
\usepackage{bm}
\hypersetup{colorlinks=true,linkcolor=blue,citecolor=blue,urlcolor=blue}

\newtheorem{theorem}{Theorem}[section]
\newtheorem{lemma}[theorem]{Lemma}
\newtheorem{corollary}[theorem]{Corollary}
\newtheorem{proposition}[theorem]{Proposition}
\theoremstyle{definition}
\newtheorem{definition}{Definition}[section]
\theoremstyle{remark}
\newtheorem{remark}{Remark}[section]

\newcommand{\s}{\sigma}
\newcommand{\e}{\varepsilon}
\newcommand{\lQ}{\ell_Q}
\newcommand{\mQ}{m_Q}
\newcommand{\Mpl}{M_{\mathrm{Pl}}}
\newcommand{\lPl}{\ell_{\mathrm{Pl}}}
\newcommand{\half}{\tfrac{1}{2}}
\newcommand{\rd}{\mathrm{d}}
\newcommand{\rs}{r_{\mathrm{s}}}
\newcommand{\nn}{\nonumber}
\newcommand{\be}{\begin{equation}}
\newcommand{\ee}{\end{equation}}
\newcommand{\bea}{\begin{eqnarray}}
\newcommand{\eea}{\end{eqnarray}}
\newcommand{\etal}{{\it et al.}}

\begin{document}

\title{Quantum Gravitational Dynamics: Master Metric, Gauge Structure, and the Classical Limit}
\author{Romeo Matshaba}
\affiliation{University of South Africa, Pretoria, South Africa}
\date{March 2026}

\begin{abstract}
We present the complete geometric foundation of Quantum Gravitational Dynamics (QGD). The metric is constructed algebraically from the graviton field $\sigma_\mu$ via the master metric, which includes a quantum stiffness term proportional to $\kappa\ell_Q^2$. In the limit $\ell_Q \to 0$, the theory reduces exactly to general relativity. We establish the Isotropic Gauge Mapping Theorem, proving that for a single static mass, the QGD metric transforms to the isotropic Schwarzschild form, and that the radial metric component $g_{rr}$ is a pure gauge artifact while the temporal component $\sigma_t$ encodes the physical gravitational field. Source signatures $\varepsilon_a = \pm 1$ encode attraction/repulsion for mass, charge, spin, and cosmological constant. The three-tier hierarchy $\eta_{\alpha\beta} - \sum\varepsilon_a\sigma_\alpha\sigma_\beta - \kappa\ell_Q^2\partial\sigma\partial\sigma$ provides a clean separation between flat spacetime, classical gravity, and quantum stiffness. This work establishes the foundational geometric structure upon which all subsequent QGD results are built.
\end{abstract}

\maketitle

\section{Introduction}

Quantum Gravitational Dynamics (QGD) reformulates gravity in terms of a fundamental field $\sigma_\mu$, derived from the phase of a Dirac spinor in the WKB limit. The metric is not fundamental but emerges algebraically from $\sigma_\mu$ via the master metric. This construction automatically incorporates a quantum stiffness term that regularizes short-distance behavior. In the classical limit $\ell_Q \to 0$, the theory reduces exactly to general relativity, providing a smooth connection between the quantum and classical regimes.

In this paper we lay the geometric foundations of QGD. We present the full master metric, explain the role of the tetrad $T^\alpha{}_\mu$ and scaling matrix $M_{\alpha\beta}$, and tabulate the source signatures $\varepsilon_a$ that encode the nature of each source (attractive or repulsive). We then prove the Isotropic Gauge Mapping Theorem, which establishes that for a single static mass, the QGD metric can be transformed to the isotropic Schwarzschild form, and that the radial metric component is a gauge artifact while the temporal component $\sigma_t$ carries the physical information. This theorem is essential for understanding the gauge structure of QGD and for all subsequent post-Newtonian calculations.

\section{The Master Metric}

The fundamental object in QGD is the graviton field $\sigma_\mu = \hbar/(mc)\,\partial_\mu S$, derived from the phase of a Dirac spinor in the WKB limit. The metric is constructed algebraically in a local Lorentz frame as $g_{\mu\nu}^{\text{(local)}} = \eta_{\mu\nu} - \sigma_\mu\sigma_\nu$. To obtain the metric in global coordinates, we introduce a tetrad $T^\alpha{}_\mu$ and a geometric scaling matrix $M_{\alpha\beta}$. The full master metric for $N$ sources is

\begin{equation}
\boxed{g_{\mu\nu}(x) = T^\alpha{}_\mu\,T^\beta{}_\nu\!\left(M_{\alpha\beta}\circ\!\left[\eta_{\alpha\beta} - \sum_{a=1}^N\varepsilon_a\,\sigma_\alpha^{(a)}\sigma_\beta^{(a)} - \kappa\ell_Q^2\,\partial_\alpha\sigma^\gamma\partial_\beta\sigma_\gamma\right]\right).}
\label{eq:master_metric}
\end{equation}

The symbol $\circ$ denotes element-wise (Hadamard) multiplication. The components have the following interpretation:

\begin{itemize}
\item $T^\alpha{}_\mu$ is the local tetrad that converts the local Lorentz frame to global coordinates. For example, in spherical coordinates, $T^\alpha{}_\mu = \operatorname{diag}(1,1,r,r\sin\theta)$.
\item $M_{\alpha\beta}$ is a geometric scaling matrix; for most applications $M_{\alpha\beta} = \delta_{\alpha\beta}$.
\item $\varepsilon_a = \pm 1$ is the source signature, encoding whether the source is attractive ($+1$) or repulsive ($-1$). The physical origins are summarized in Table~\ref{tab:sources}.
\item $\sigma_\mu^{(a)}$ is the graviton field of source $a$. For a point mass, $\sigma_t^{(a)} = \sqrt{2GM_a/(c^2 r_a)}$, with all spatial components zero in the rest frame.
\item $\ell_Q = \sqrt{G\hbar^2/c^4} \sim 10^{-35}$ m is the quantum length scale, and $\kappa \sim 2$ is the stiffness coefficient. The term $\kappa\ell_Q^2\,\partial_\alpha\sigma^\gamma\partial_\beta\sigma_\gamma$ represents quantum stiffness; it is negligible at macroscopic scales but essential for singularity resolution.
\end{itemize}

Equation~\eqref{eq:master_metric} exhibits a clean three-tier hierarchy:

\[
g_{\mu\nu} = \underbrace{T T(\eta)}_{\text{Tier 0: Minkowski}} - \underbrace{T T(\sum\varepsilon\sigma\sigma)}_{\text{Tier 1: Classical gravity}} - \underbrace{T T(\kappa\ell_Q^2\partial\sigma\partial\sigma)}_{\text{Tier 2: Quantum stiffness}}.
\]

In the classical limit $\ell_Q \to 0$, the quantum stiffness term vanishes, and the metric reduces to the standard GR form built from the $\sigma$-fields. This limit defines the classical equivalence between QGD and general relativity.

\begin{table}[h]
\centering
\caption{Source signatures in QGD. Each source contributes a $\sigma$-field with a definite sign $\varepsilon_a$, encoding whether the interaction is attractive ($+1$) or repulsive ($-1$).}
\label{tab:sources}
\begin{tabular}{llll}
\toprule
Source & $\sigma$-field & $\varepsilon_a$ & Effect \\
\midrule
Mass & $\displaystyle\sqrt{\frac{2GM}{c^2 r}}$ & $+1$ & Attraction \\
Spin (Kerr) & $\displaystyle a\sin\theta\sqrt{\frac{2GM}{c^2 r\,\Sigma}}$ & $+1$ & Frame dragging \\
Electric charge & $\displaystyle\sqrt{\frac{GQ^2}{c^4 r^2}}$ & $-1$ & Repulsion \\
Cosmological constant ($\Lambda>0$) & $\displaystyle Hr/c$ & $-1$ & de Sitter expansion \\
Cosmological constant ($\Lambda<0$) & $\displaystyle Hr/c$ & $+1$ & Anti-de Sitter confinement \\
\bottomrule
\end{tabular}
\end{table}

\section{The Isotropic Gauge and Radial Gauge Artifact}

In QGD, we have the freedom to choose a gauge for the spatial part of the $\sigma$-field. The isotropic gauge is defined by setting $\sigma_r = 0$ for all sources. In this gauge, the spatial metric is flat: $g_{ij} = \delta_{ij}$. However, this does not mean spacetime is flat; the curvature is encoded in $\sigma_t$ and its gradients. The choice $\sigma_r = 0$ is analogous to the Coulomb gauge in electromagnetism, where the vector potential is transverse and the scalar potential carries the physical information.

We now prove that the radial metric component $g_{rr}$ is a pure gauge artifact. Consider a single static mass. In isotropic gauge, the QGD metric is

\[
ds^2 = -(1 - \sigma_t^2) dt^2 + dr^2 + r^2 d\Omega^2,
\]

with $\sigma_t = \sqrt{2GM/(c^2 r)}$. Perform a coordinate transformation to a new radial coordinate $\rho$ defined by

\[
r = \rho\left(1 + \frac{GM}{2c^2\rho}\right)^2.
\]

A straightforward calculation yields

\[
ds^2 = -\left(\frac{1 - \frac{GM}{2c^2\rho}}{1 + \frac{GM}{2c^2\rho}}\right)^2 dt^2 + \left(1 + \frac{GM}{2c^2\rho}\right)^4 (d\rho^2 + \rho^2 d\Omega^2),
\]

which is exactly the isotropic Schwarzschild metric. In these coordinates, $g_{\rho\rho} \neq 1$, but the physical content is unchanged. The graviton scalar transforms covariantly and becomes $\sigma_t = \sqrt{2GM/(c^2\rho)}$ in the new coordinates. Thus $\sigma_t$ is invariant under spatial coordinate changes (it is a scalar), while $g_{rr}$ changes. This proves that $g_{rr}$ is a gauge artifact; the physical gravitational field is encoded entirely in $\sigma_t$.

\section{The Isotropic Gauge Mapping Theorem}

We now state and prove the central theorem that establishes the equivalence between QGD and GR in the classical limit.

\begin{theorem}[Isotropic Gauge Mapping]
\label{thm:isotropic}
Let $g_{\mu\nu}$ be the QGD master metric for a single static mass $M$ in the classical limit $\ell_Q \to 0$, with $\sigma_t = \sqrt{2GM/(c^2 r)}$ and $\sigma_i = 0$ in isotropic gauge. Then there exists a coordinate transformation to a new radial coordinate $\rho$ such that the metric takes the exact isotropic Schwarzschild form
\begin{equation}
ds^2 = -\left(\frac{1 - \frac{GM}{2c^2\rho}}{1 + \frac{GM}{2c^2\rho}}\right)^2 c^2 dt^2 + \left(1 + \frac{GM}{2c^2\rho}\right)^4 (d\rho^2 + \rho^2 d\Omega^2).
\end{equation}
Moreover, the graviton scalar $\sigma_t$ is invariant under this transformation, and the radial metric component $g_{rr}$ is a pure gauge artifact.
\end{theorem}

\begin{proof}
In isotropic gauge, the QGD metric is
\begin{equation}
ds^2 = -(1 - \sigma_t^2) dt^2 + dr^2 + r^2 d\Omega^2, \qquad \sigma_t = \sqrt{\frac{2GM}{c^2 r}}.
\end{equation}
Define a new radial coordinate $\rho$ by the relation
\begin{equation}
r = \rho\left(1 + \frac{GM}{2c^2\rho}\right)^2.
\end{equation}
This transformation is invertible for $\rho > GM/(2c^2)$. The differential is
\begin{equation}
dr = \left(1 + \frac{GM}{2c^2\rho}\right)^2 d\rho - 2\rho\left(1 + \frac{GM}{2c^2\rho}\right)\frac{GM}{2c^2\rho^2} d\rho = \left(1 - \frac{GM}{2c^2\rho}\right)\left(1 + \frac{GM}{2c^2\rho}\right) d\rho.
\end{equation}
Hence
\begin{equation}
dr^2 = \left(1 - \frac{GM}{2c^2\rho}\right)^2\left(1 + \frac{GM}{2c^2\rho}\right)^2 d\rho^2.
\end{equation}
Also,
\begin{equation}
r^2 = \rho^2\left(1 + \frac{GM}{2c^2\rho}\right)^4.
\end{equation}
Now express $\sigma_t$ in terms of $\rho$. From the definition,
\begin{equation}
\sigma_t = \sqrt{\frac{2GM}{c^2 r}} = \sqrt{\frac{2GM}{c^2\rho}} \left(1 + \frac{GM}{2c^2\rho}\right)^{-1}.
\end{equation}
Therefore,
\begin{align}
1 - \sigma_t^2 &= 1 - \frac{2GM}{c^2\rho}\left(1 + \frac{GM}{2c^2\rho}\right)^{-2} \nn\\
&= \frac{\left(1 + \frac{GM}{2c^2\rho}\right)^2 - \frac{2GM}{c^2\rho}}{\left(1 + \frac{GM}{2c^2\rho}\right)^2} = \frac{\left(1 - \frac{GM}{2c^2\rho}\right)^2}{\left(1 + \frac{GM}{2c^2\rho}\right)^2}.
\end{align}
Substituting into the line element,
\begin{align}
ds^2 &= -\frac{\left(1 - \frac{GM}{2c^2\rho}\right)^2}{\left(1 + \frac{GM}{2c^2\rho}\right)^2} dt^2 + \left(1 - \frac{GM}{2c^2\rho}\right)^2\left(1 + \frac{GM}{2c^2\rho}\right)^2 d\rho^2 \nn\\
&\quad + \rho^2\left(1 + \frac{GM}{2c^2\rho}\right)^4 d\Omega^2 \nn\\
&= -\left(\frac{1 - \frac{GM}{2c^2\rho}}{1 + \frac{GM}{2c^2\rho}}\right)^2 dt^2 + \left(1 + \frac{GM}{2c^2\rho}\right)^4 \left(d\rho^2 + \rho^2 d\Omega^2\right).
\end{align}
This is exactly the isotropic Schwarzschild metric. The graviton scalar transforms to $\sigma_t = \sqrt{2GM/(c^2\rho)}$, which is the same functional form, demonstrating its invariance under the transformation. The radial metric component changes from $g_{rr}=1$ to $g_{\rho\rho} = (1 + GM/(2c^2\rho))^4$, proving that $g_{rr}$ is a gauge artifact. ∎
\end{proof}

\section{Equivalence to General Relativity in the Classical Limit}

Theorem~\ref{thm:isotropic} establishes that for a single static mass, the QGD metric can be mapped to the isotropic Schwarzschild metric. Since the isotropic Schwarzschild metric is a solution of Einstein's equations, and Birkhoff's theorem guarantees uniqueness, any spherically symmetric vacuum solution of Einstein's equations must be isometric to this metric. Therefore, the QGD master metric with $\ell_Q = 0$ produces exactly the same set of spherically symmetric vacuum solutions as general relativity.

For the two-body problem, the situation is more subtle, but the same principle applies: the QGD field equations are derived from the action, and in the limit $\ell_Q \to 0$, they reduce to the Einstein equations (as shown by the equivalence theorem in Chapter~3). The Isotropic Gauge Mapping Theorem for the single-body case serves as a foundational consistency check, confirming that QGD reproduces GR in the appropriate limit.

\section{Conclusion}

We have presented the complete geometric foundation of QGD. The master metric, Eq.~\eqref{eq:master_metric}, provides an algebraic construction of spacetime from the graviton field $\sigma_\mu$, with a three-tier hierarchy separating flat spacetime, classical gravity, and quantum stiffness. Source signatures $\varepsilon_a = \pm 1$ encode attraction and repulsion for mass, spin, charge, and cosmological constant. The Isotropic Gauge Mapping Theorem proves that for a single static mass, the QGD metric can be transformed to the isotropic Schwarzschild form, establishing that $g_{rr}$ is a pure gauge artifact while $\sigma_t$ carries the physical gravitational field. In the classical limit $\ell_Q \to 0$, QGD reduces exactly to general relativity.

These results lay the groundwork for all subsequent developments in QGD, including the post-Newtonian expansion, the derivation of the EOB potential, and the study of gravitational waves and black hole thermodynamics. The gauge structure revealed here is essential for understanding the role of $\sigma_r$ as a gauge field and the physical significance of $\sigma_t$.

\end{document}
